\subsection{实验方案与评价协议}
\label{sec:task1-protocol}
表~\ref{tab:task1-design} 列出预先固定的四个方案。A 与 B 检验整体表示的差异，其中预处理分别为整图缩放与中心裁剪；因此它们并不是只改变网络结构的严格单因素对照。B 与 C 检验 PCA 压缩，B 与 D 则在完全相同的特征空间比较聚类算法。

\begin{table*}[t]
\caption{预先固定的实验方案。所有方法均将 600 张图像划分为六簇。}
\label{tab:task1-design}
\centering
\begin{tabular}{clcl}
\toprule
编号 & 输入表示 & 维度 & 聚类算法\\
\midrule
A & HSV + uniform LBP，直方图平方根变换及 $L_2$ 归一化 & 138 & K-means，五个固定种子\\
B（主方法） & 冻结 ResNet-18，全局平均池化后 $L_2$ 归一化 & 512 & K-means，五个固定种子\\
C & B 的特征经 PCA 保留 64 维后再次 $L_2$ 归一化 & 64 & K-means，五个固定种子\\
D & 与 B 相同的归一化深度特征 & 512 & Ward 凝聚聚类\\
\bottomrule
\end{tabular}
\end{table*}

\subsubsection{外部评价：与类别标签的一致性}
令 $n_{ab}$ 表示真实类别为 $a$、预测簇为 $b$ 的样本数。聚类准确率 ACC 先求真实类别与簇编号的一对一最优匹配：
\begin{equation}
\mathrm{ACC}=\frac{1}{N}\max_{\pi}\sum_{i=1}^{N}
\ind[y_i=\pi(c_i)].
\end{equation}
其中 $\pi$ 为六个编号的置换，通过线性指派求解。该映射只用于评价和显示混淆矩阵，不返回训练过程，也不等于已经学会为新图像赋予语义标签。

调整兰德指数 ARI\citep{hubert1985} 评价两组划分对样本对的同簇判断一致性，并扣除随机一致性的期望。令 $a_a=\sum_b n_{ab}$、$b_b=\sum_a n_{ab}$，定义
\begin{align}
T&=\sum_{ab}\binom{n_{ab}}{2},\\
A&=\sum_a\binom{a_a}{2},\quad
B=\sum_b\binom{b_b}{2},\quad
E=\frac{AB}{\binom{N}{2}},\\
\mathrm{ARI}&=\frac{T-E}{(A+B)/2-E}.
\end{align}
ARI 为 1 表示两个划分完全一致，随机划分的期望接近 0，负值表示低于相应随机基准。ARI 对簇编号置换不敏感。

归一化互信息 NMI 使用算术平均熵归一化：
\begin{equation}
\mathrm{NMI}=\frac{2I(Y;C)}{H(Y)+H(C)}.
\end{equation}
$H$ 为经验熵，$I$ 为真实类别与预测簇的互信息。NMI 位于 $[0,1]$，较高值表示两个划分共享较多信息，但它不作 ARI 式的随机一致性校正。三个外部指标统一报告，避免单个指标遗漏簇合并、簇拆分或类别不平衡的影响。

\subsubsection{内部评价：特征空间中的簇结构}
内部评价不使用真实标签。对每个样本，令 $a(i)$ 为同簇平均距离，$b(i)$ 为与其他各簇平均距离的最小值，轮廓系数为
\begin{equation}
s(i)=\frac{b(i)-a(i)}{\max\{a(i),b(i)\}},\qquad
S=\frac{1}{N}\sum_i s(i).
\end{equation}
计算采用欧氏距离和全部 600 张图像；若出现单样本簇，则按实现约定令该样本轮廓系数为零。较大的 $S$ 表示相对更紧凑、分离的几何结构，并不自动等价于更准确的语义类别。

Calinski--Harabasz（CH）指标比较簇间与簇内离散程度：
\begin{equation}
\mathrm{CH}=\frac{\sum_k n_k\|\mu_k-\bar z\|_2^2/(K-1)}
{\sum_k\sum_{i:c_i=k}\|z_i-\mu_k\|_2^2/(N-K)}.
\end{equation}
Davies--Bouldin（DB）指标采用簇内平均半径
$S_k=n_k^{-1}\sum_{i:c_i=k}\|z_i-\mu_k\|_2$：
\begin{equation}
\mathrm{DB}=\frac{1}{K}\sum_k\max_{\ell\ne k}
\frac{S_k+S_\ell}{\|\mu_k-\mu_\ell\|_2}.
\end{equation}
CH 越高、DB 越低通常表示更清晰的紧凑簇结构。具体计算与边界情况采用 scikit-learn 的标准实现\citep{sklearnclustering}。

每个方案在其自身输入特征上计算内部指标。A、B、C 使用不同维度和不同几何表示，因此不能仅凭其轮廓系数、CH 或 DB 数值大小直接判定跨表示的语义性能。B 与 D 的内部指标具有相同特征和距离基础，可直接用于这两个算法的结构比较。

\subsubsection{随机稳定性与可核查性}
五个固定种子分别报告外部、内部指标，并汇总均值与标准差。另计算五个输出划分之间共 $\binom52=10$ 对 ARI，描述初始化改变时簇划分的一致程度；这里两组输入均为算法输出，完全不使用类别标签。高稳定性说明结果对这些初始化较一致，不保证划分语义正确，也不表示对新样本或数据扰动具有稳定性。

主结果的混淆矩阵按最优标签匹配重排，用于定位类别混淆；二维投影仅辅助观察，不代替在原聚类特征上计算指标。每簇样例用于解释外观相似性与错误归并，不能仅凭少量展示样本推断全部样本表现。报告中的数值以保存的逐次运行结果为准，保持模型配置、随机种子和评价口径一致。
